\documentclass[../../../include/open-logic-section]{subfiles}

\begin{document}

\olfileid[fr]{his}{set}{mythology}

\olsection{Plus de mythe que d'histoire ?}

Avec plus d'un siècle de recul, nous devons nous garder de reprendre
ces jugements sans examen. Ces résultats étaient assurément nouveaux,
stimulants et surprenants. Mais étaient-ils vraiment aussi
bouleversants ? Montraient-ils réellement qu'il fallait renoncer à
l'intuition géométrique ?

Sur le premier point, \citet{Gouvea2011} observe que les célèbres
mots de Cantor à Dedekind, «~\emph{je le vois, mais je ne le crois
pas}~», sont souvent cités hors de leur contexte. En voici un
extrait plus long, cité par \citeauthor{Gouvea2011} :
\begin{quote}
Veuillez excuser l'ardeur avec laquelle je m'intéresse à ce sujet,
si je sollicite si souvent votre bonté et votre patience. Les
communications que je vous ai envoyées récemment sont, même pour
moi, si inattendues, si nouvelles, que je ne trouverai pas la
tranquillité d'esprit tant que je n'aurai pas reçu de vous, cher ami,
un jugement sur leur exactitude. Tant que vous ne m'aurez pas donné votre
accord, je ne peux que dire : \emph{je le vois, mais je ne le crois
pas.}
\end{quote}
Cantor savait que son résultat était «~si inattendu, si nouveau~».
Il est cependant douteux qu'il l'ait jamais trouvé
\emph{incroyable}. Comme le souligne \citeauthor{Gouvea2011}, il
demandait simplement à Dedekind de vérifier la démonstration qu'il
lui avait soumise.

Quant à l'intuition géométrique, Peano publia sa courbe remplissant
l'espace sans aucune figure. Mais lorsque Hilbert publia la sienne,
il expliqua son objectif : offrir à ses lecteurs un moyen clair de
comprendre le résultat de Peano, pourvu qu'ils «~s'aident de
l'intuition géométrique suivante~». Il présenta alors une série de
\emph{figures} semblables à celles de la
\olref[his][set][pathology]{sec}.

Plus généralement, même si les figures sont passées quelque peu de
mode dans les démonstrations publiées, les mathématiciens les
utilisent \emph{fréquemment} pour démontrer des résultats. En
première approximation, les bons mathématiciens savent quand ils
peuvent se fier à l'intuition géométrique.

Bref, ne cédons pas à l'emballement, ou du moins ne reprenons pas
ce récit sans examen. Pour approfondir la question,
on pourra lire \citet{Giaquinto2007}.

\end{document}
